\begin{theorem}[Affine pencil of quadratic phases]\label{mod:finitefield-affinepencil}
Let $p$ be an odd prime, let $k$ be a finite field of characteristic
$p$, and write $[j]\in k$ for the canonical image of
$j\in\F_p=\mathbf Z/p\mathbf Z$.  Let $\psi:k\to\C^\times$ be a
nontrivial additive character satisfying
\[
  \psi(z^p)=\psi(z)\qquad(z\in k).
\]
Suppose that $\rho,\sigma,\tau,a_0,b_0\in k$ satisfy
\[
 \rho\ne0,\qquad 1+[j]\rho\ne0\quad(j\in\F_p),
\]
and
\[
 a_0^p=1-(\rho^{p-1})^{-1},
 \qquad b_0^p=\sigma-\rho^{-1}\tau.
\]
Then, for the normalized quadratic phase of
Theorem~\ref{mod:finitefield-quadraticphase},
\[
 Q_\psi(a_0,b_0)
 \prod_{j\in\F_p^\times}Q_\psi([j]\rho,[j]\tau)
 =\prod_{j\in\F_p}Q_\psi(1+[j]\rho,\sigma+[j]\tau).
\]

In particular, all quadratic coefficients occurring here are nonzero;
this is a consequence of the displayed hypotheses, rather than an
additional assumption.  The proof includes the exact rational identity
\[
 \sum_{j\in\F_p}\frac{(\sigma+[j]\tau)^2}{1+[j]\rho}
 =\frac{\rho^{p-3}(\sigma\rho-\tau)^2}{\rho^{p-1}-1},
\]
with the denominator proved nonzero from
$\prod_j(1+[j]\rho)=1-\rho^{p-1}$.

The identity is homogeneous in the common coefficient scale in the precise
form required at the wild odd boundary.  Namely, let $c,c_0\in k$ satisfy
$c\ne0$ and $c_0^p=c$.  Then
\[
 \frac{
   Q_\psi(ca_0^p,cb_0^p)
   \displaystyle\prod_{u\in\F_p^\times}
     Q_\psi(c[u]\rho,c[u]\tau)}{
   \displaystyle\prod_{j\in\F_p}
     Q_\psi(c(1+[j]\rho),c(\sigma+[j]\tau))}
 =
 \frac{
   Q_\psi(a_0,b_0)
   \displaystyle\prod_{u\in\F_p^\times}
     Q_\psi([u]\rho,[u]\tau)}{
   \displaystyle\prod_{j\in\F_p}
     Q_\psi(1+[j]\rho,\sigma+[j]\tau)}.
\]
Indeed, with $z_0=-b_0^2/(2a_0)$, the nonzero norm rows contribute
$\nu(c)^{p-1}=1$ and no additive factor because
$\sum_{u\in\F_p^\times}[u]=0$.  The affine rational-sum identity makes the
twist-row exponent equal to $z_0^p$.  Frobenius invariance and $c_0^p=c$
give
\[
 \nu(c)^p=\nu(c_0),\qquad
 \psi((c-1)z_0^p)=\psi((c_0-1)z_0).
\]
Thus the numerator and denominator acquire the same explicitly nonzero
factor, which is cancelled only after its nonvanishing is proved.
\LeanFile{LanglandsFirstMainLemma/FiniteField/AffinePencil.lean}
\lean{LanglandsFirstMainLemma.affinePencil}
\lean{LanglandsFirstMainLemma.affinePencil_homogeneous_quotient}
\leanok
\SourceText{Lemma lem:affine-pencil and the boundary case of Proposition prop:odd-phase-cancel.}
\uses{mod:finitefield-quadraticphase,mod:basic-finiteproducts}
\end{theorem}
